\section*{AGRADECIMIENTOS}

Quiero agradecer a la Universidad San Francisco de Quito y al Instituto de Microbiología por la oportunidad de estudiar en este programa de maestría mediante una beca, sin la cual, alcanzar esta meta no habría sido posible.

Expreso también mi más profundo agradecimiento hacia mi tutor, el doctor Paúl Cárdenas, por la orientación brindada durante el desarrollo de este trabajo y por la confianza depositada en mis capacidades. Sus observaciones y recomendaciones contribuyeron a orientar el proyecto y a superar los diferentes desafíos que surgieron durante el proceso. 

A los miembros de mi comité, Verónica Barragán, Shazia Ruybal y Henry Vaca, les agradezco por su compromiso, retroalimentación valiosa y por estar pendientes del avance de este proyecto.

Al equipo de Bioinformática, especialmente a Erika Muñoz y Rommel Guevara, gracias por sus enseñanzas y paciencia para el desarrollo de mi proyecto, su ayuda fue fundamental para ejecutar este trabajo.

A la Planta de Tratamiento de Aguas Residuales de Quitumbe, sobre todo al equipo de la Unidad de operaciones, quienes con la mejor actitud, siempre estuvieron dispuestos a ayudarme, especialmente a Juan Guadamud, gracias por la colaboración con la obtención de las muestras semanales. 

A mis amigos de la cohorte 2024: Anahí, Edison, Melissa, Jerram, Alex, David y Andrea, muchas gracias por hacer de esta etapa una experiencia tan especial, me llevo en el corazón todos los momentos que compartimos. 

A mis hermanas, Lissette y Doménica por saber cómo sacarme una sonrisa incluso en los días más difíciles y por escucharme cuando más lo necesitaba. Gracias por su cariño y por hacer que este camino fuera un poco más llevadero.

Finalmente, quiero agradecer a mis padres, Ofelia Sánchez y Edwin Viscarra, que siempre han estado para mí en cada paso de mi vida. Sin su amor y apoyo, este logro no habría sido posible. 

